\section{The qutrit family for entropy scaling}

The following family is the state introduced in~\cite{Schuch2008EntropyScaling},
arXiv:0705.0292v2, at the start of the examples.

\begin{definition}[Qutrit state and its weights]
    \label{def:qutrit_entropy_scaling_state}
    \lean{TNLean.MPS.AreaLaw.EntropyScaling.QutritLabel,
        TNLean.MPS.AreaLaw.EntropyScaling.qutritWord,
        TNLean.MPS.AreaLaw.EntropyScaling.qutritWeight,
        TNLean.MPS.AreaLaw.EntropyScaling.qutritState,
        TNLean.MPS.AreaLaw.EntropyScaling.qutritFamily,
        TNLean.MPS.AreaLaw.EntropyScaling.qutritHalves,
        TNLean.MPS.AreaLaw.EntropyScaling.qutritHalves_append}
    \leanok
    Let $N$ be a positive integer and let $0\le p\le1$. On $2N$ qutrits define
    \begin{align}
        \ket{\psi_{2N}(p)}
        =\sqrt{1-p}\ket{2}^{\otimes2N}
        +\sqrt{p/2^N}\sum_{x\in{\{0,1\}}^N}\ket{x}\ket{x}.
        \label{eq:qutrit_entropy_scaling_state}
    \end{align}
    The coefficient labels are a flag $\star$ and the binary words $x$ of
    length $N$. Their half-chain words are $w_\star=(2,\ldots,2)$ and $w_x=x$,
    and their squared coefficients are
    $\lambda_\star=1-p$ and $\lambda_x=p/2^N$.
    Splitting a chain configuration after its first $N$ sites gives its two
    half-chain configurations. The source's family takes $p=p_N=1/N$.
\end{definition}

\begin{theorem}[Orthonormal Schmidt vectors]
    \label{thm:qutrit_entropy_scaling_orthonormal}
    \lean{TNLean.MPS.AreaLaw.EntropyScaling.orthonormal_qutritWord,
        TNLean.MPS.AreaLaw.EntropyScaling.qutritWord_injective,
        TNLean.MPS.AreaLaw.EntropyScaling.orthonormal_qutritFullWord,
        TNLean.MPS.AreaLaw.EntropyScaling.qutritBasisMatrix,
        TNLean.MPS.AreaLaw.EntropyScaling.qutritBasisMatrix_conjTranspose_mul_self,
        TNLean.MPS.AreaLaw.EntropyScaling.qutritBasisMatrix_conjTranspose}
    \leanok
    \uses{def:qutrit_entropy_scaling_state}
    The vectors $e_i=\ket{w_i}$ are a real orthonormal family, as are
    $\ket{w_i}\ket{w_i}$. Thus the matrix $B$ with columns $e_i$ satisfies
    $B^\dagger B=1$ and $B^\dagger=B^{\mathsf T}$.
\end{theorem}
\begin{proof}\leanok
    The binary words are distinct, and no binary word equals the flag word
    when $N>0$. Apply orthonormality of the computational basis and of its
    doubled configurations.
\end{proof}

\begin{theorem}[Normalization]
    \label{thm:qutrit_entropy_scaling_normalization}
    \lean{TNLean.MPS.AreaLaw.EntropyScaling.norm_qutritState,
        TNLean.MPS.AreaLaw.EntropyScaling.qutritWeight_nonneg,
        TNLean.MPS.AreaLaw.EntropyScaling.sum_qutritWeight,
        TNLean.MPS.AreaLaw.EntropyScaling.norm_qutritFamily}
    \leanok
    \uses{def:qutrit_entropy_scaling_state}
    The weights $\lambda_i$ are nonnegative and sum to one, and
    $\|\psi_{2N}(p)\|=1$. In particular, $\|\psi_{2N}(1/N)\|=1$.
\end{theorem}
\begin{proof}\leanok
    \uses{thm:qutrit_entropy_scaling_orthonormal}
    There are $2^N$ binary words, so
    $\sum_i\lambda_i=1-p+2^N(p/2^N)=1$.
    Orthonormality in~(\ref{eq:qutrit_entropy_scaling_state}) gives the
    asserted norm. For $N\ge1$, one has $0<1/N\le1$.
\end{proof}

\begin{definition}[Middle coefficient matrix and marginal]
    \label{def:qutrit_entropy_scaling_middle}
    \lean{TNLean.MPS.AreaLaw.EntropyScaling.qutritMiddleVector,
        TNLean.MPS.AreaLaw.EntropyScaling.qutritMiddleMatrix,
        TNLean.MPS.AreaLaw.EntropyScaling.qutritMiddleMarginal}
    \leanok
    \uses{def:qutrit_entropy_scaling_state}
    Write the state in the two half-chain coordinates, and let
    $C_{a,b}=\langle a,b\mid\psi_{2N}(p)\rangle$.
    Its left marginal is
    $\rho=\operatorname{Tr}_{\mathrm{right}}
        \ket{\psi_{2N}(p)}\!\bra{\psi_{2N}(p)}=CC^\dagger$.
\end{definition}

\begin{theorem}[Identification with the chain marginal]
    \label{thm:qutrit_entropy_scaling_native_marginal}
    \lean{TNLean.MPS.AreaLaw.EntropyScaling.qutritLeftConfigurationEquiv,
        TNLean.MPS.AreaLaw.EntropyScaling.qutritRightConfigurationEquiv,
        TNLean.MPS.AreaLaw.EntropyScaling.qutritConfigurations_cutEquiv,
        TNLean.MPS.AreaLaw.EntropyScaling.cutMatrix_qutritState_reindex,
        TNLean.MPS.AreaLaw.EntropyScaling.regionState_qutritState_reindex}
    \leanok
    \uses{def:qutrit_entropy_scaling_middle}
    Numbering the sites in each half from zero identifies $C$ with the
    coefficient matrix of the actual chain cut after its first $N$ sites.
    Under the same relabelling on the left half, $\rho$ is the regional
    state of these first $N$ sites.
\end{theorem}
\begin{proof}\leanok
    Restrict a configuration to the first half and to the second half,
    subtracting $N$ from the site indices of the latter. This is the same
    decomposition as the two half-chain coordinates. Relabel the cut matrix
    accordingly. Its product with its adjoint gives the asserted marginal.
\end{proof}

\begin{theorem}[Middle-cut Schmidt factorization]
    \label{thm:qutrit_entropy_scaling_schmidt}
    \lean{TNLean.MPS.AreaLaw.EntropyScaling.qutritMiddleMatrix_factorization,
        TNLean.MPS.AreaLaw.EntropyScaling.qutritMiddleVector_eq}
    \leanok
    \uses{def:qutrit_entropy_scaling_middle,
        thm:qutrit_entropy_scaling_orthonormal}
    With $B$ as above, $C=B\operatorname{diag}(\sqrt{\lambda_i})B^\dagger$.
\end{theorem}
\begin{proof}\leanok
    Split each doubled word in~(\ref{eq:qutrit_entropy_scaling_state}) at
    the middle cut. Its two factors are both $e_i$, with coefficient
    $\sqrt{\lambda_i}$. Since the vectors $e_i$ are real, their transpose
    equals their adjoint in the coefficient matrix.
\end{proof}

\begin{theorem}[Diagonalization of the middle marginal]
    \label{thm:qutrit_entropy_scaling_marginal}
    \lean{TNLean.MPS.AreaLaw.EntropyScaling.qutritMiddleMarginal_factorization}
    \leanok
    \uses{def:qutrit_entropy_scaling_middle,
        thm:qutrit_entropy_scaling_orthonormal}
    The left marginal satisfies $\rho=B\operatorname{diag}(\lambda_i)B^\dagger$.
\end{theorem}
\begin{proof}\leanok
    \uses{thm:qutrit_entropy_scaling_schmidt,
        thm:qutrit_entropy_scaling_orthonormal}
    Substitute the Schmidt factorization into $CC^\dagger$, use
    $B^\dagger B=1$, and square the nonnegative square roots of the weights.
\end{proof}

\begin{theorem}[Complete middle-cut spectrum]
    \label{thm:qutrit_entropy_scaling_spectrum}
    \lean{TNLean.MPS.AreaLaw.EntropyScaling.charpoly_qutritMiddleMarginal,
        TNLean.MPS.AreaLaw.EntropyScaling.charpoly_qutritFamily_middle,
        TNLean.MPS.AreaLaw.EntropyScaling.charpoly_regionState_qutritState}
    \leanok
    \uses{def:qutrit_entropy_scaling_middle}
    The characteristic polynomial of the middle marginal, and therefore of
    the actual first-$N$-site regional state, is
    \begin{align}
        \det(tI-\rho)
        =t^{3^N-1-2^N}\bigl(t-(1-p)\bigr)
        {\bigl(t-p/2^N\bigr)}^{2^N}.
        \label{eq:qutrit_entropy_scaling_charpoly}
    \end{align}
    Thus its squared Schmidt coefficients are $1-p$, one copy of $p/2^N$
    for each binary word, and the remaining zeros. This statement includes
    zero weights at $p=0$ and $p=1$, and applies to $p=1/N$.
\end{theorem}
\begin{proof}\leanok
    \uses{thm:qutrit_entropy_scaling_marginal,
        thm:qutrit_entropy_scaling_native_marginal,
        thm:qutrit_entropy_scaling_orthonormal}
    The injective map from the labels to the half-chain words gives
    $1+2^N\le3^N$. Apply the rectangular determinant identity to
    $B\operatorname{diag}(\lambda_i)B^\dagger$ and
    $\operatorname{diag}(\lambda_i)B^\dagger B
        =\operatorname{diag}(\lambda_i)$. The remaining $3^N-1-2^N$
    dimensions contribute zero eigenvalues. The diagonal entries give
    the flag factor and the $2^N$ equal binary factors in
    (\ref{eq:qutrit_entropy_scaling_charpoly}).
\end{proof}
